% sec-08-previous-human-experience.tex - Previous Human Experience (full stage)
\section{Previous Human Experience}

This section assembles the previous human experience that supports the proposed
Phase 1 investigation, organized under 21 CFR \S 312.23(a)(9). It states the
marketed status of daraxonrasib (RMC-6236), summarizes the prior and ongoing
clinical research experience with the drug, situates the proposed combination
within the contemporary clinical care experience for robotic
pancreaticoduodenectomy, and identifies the safety-reporting framework under
which post-submission human experience will continue to be characterized. Because
this is a first-in-human investigation of the drug in a perioperative,
curative-intent, device-directed surgical context, the evidence base that
substitutes for direct prior experience in that exact setting is drawn from three
sources: the late-phase metastatic-disease program for daraxonrasib, the
established human experience with robotic pancreaticoduodenectomy benchmarked by
the 2025 Dutch nationwide cohort, and the author's 2025 to 2026 computational
program that supplies the in silico scaffolding for the device-directed component
\cite{daraxwhipple, onpremwhippl}. Figure 22 maps this evidence base, and the
narrative below develops each contributing strand in turn, consistent with the
phased-development expectations for an initial Phase 1 submission
\cite{phase1guidance}.

\begin{mermaidfig}
\node[mmgoal] (evid) at (0,0) {Evidence base for\\the IND (\S 3.2, \S 9)};
% Clinical lane (left)
\node[mmlanetitle] at (-5.0,2.55) {Clinical (daraxonrasib)};
\node[mmproc] (r302) at (-5.0,1.9)  {RASolute 302\\Phase 3, pretreated\\metastatic PDAC};
\node[mmproc] (r301) at (-5.0,0.0)  {RASolve 301\\first-line program};
\node[mmproc] (btd)  at (-5.0,-1.9) {Breakthrough Therapy\\(June 2025)};
% Computational lane (right)
\node[mmlanetitle] at (5.0,3.45) {Computational (author 2025-2026)};
\node[mmin] (c1) at (5.0,2.8)   {60-sec PDAC Whipple +\\daraxonrasib simulation};
\node[mmin] (c2) at (5.0,1.4)   {AI digital-twin\\PDAC simulation};
\node[mmin] (c3) at (5.0,0.0)   {QSP metastatic-PDAC\\sim (VVUQ)};
\node[mmin] (c4) at (5.0,-1.4)  {100,000-patient in silico\\Phase 3 (triplicate)};
\node[mmin] (c5) at (5.0,-2.8)  {End-to-end PDAC\\digital-twin proposals};
% Outcome
\node[mmdark] (new) at (0,-3.6) {First-in-human in the\\perioperative, curative-intent,\\device-directed context};
\draw[mmedge] (r302) to[out=0,in=150,looseness=0.7] (evid.west);
\draw[mmedge] (r301.east) to[out=0,in=180,looseness=0.6] (evid.west);
\draw[mmedge] (btd) to[out=0,in=210,looseness=0.7] (evid.west);
\draw[mmedge] (c1) to[out=180,in=30,looseness=0.7] (evid.east);
\draw[mmedge] (c2.west) to[out=180,in=20,looseness=0.6] (evid.east);
\draw[mmedge] (c3.west) to[out=180,in=0,looseness=0.6] (evid.east);
\draw[mmedge] (c4.west) to[out=180,in=-20,looseness=0.6] (evid.east);
\draw[mmedge] (c5) to[out=180,in=-30,looseness=0.7] (evid.east);
\draw[mmedge] (evid) -- (new);
\end{mermaidfig}
\vspace{-0.4cm}
\figcaption{Figure 22. The previous human experience and evidence base. Clinically,
daraxonrasib is supported by the RASolute 302 Phase 3 program in previously treated
metastatic PDAC, the RASolve 301 first-line program, and its June 2025 Breakthrough
Therapy designation; computationally, the author's 2025 to 2026 works supply the
simulation scaffolding (the 60-second PDAC Whipple and daraxonrasib simulation, the
AI digital-twin PDAC simulation, the quantitative-systems-pharmacology metastatic
PDAC simulation with verification, validation, and uncertainty quantification, the
100,000-patient in silico Phase 3 triplicate, and the end-to-end PDAC digital-twin
proposals). Daraxonrasib remains first-in-human in this perioperative,
curative-intent, device-directed surgical context.}

\subsection{Marketed experience}

Daraxonrasib (RMC-6236) is an investigational new drug and is not approved for
marketing in the United States or in any other country for any indication. It is
not available commercially, is not the subject of any New Drug Application
approval, and therefore has no marketed experience, no postmarketing surveillance
history, and no labeling derived from marketing authorization. All human exposure
to date has occurred within controlled clinical investigations conducted under
investigational status, and the safety and efficacy information available for the
drug is accordingly derived entirely from that investigational program rather than
from any marketed use \cite{daraxwhipple}. Because there is no marketed experience,
this subsection records that fact and points instead to the investigational
program characterized in \S 9.2; no withdrawal, suspension, or restriction of any
marketing authorization has occurred in any country, because none exists.

Daraxonrasib is an orally administered, once-daily RAS(ON) multi-selective
(pan-RAS) small-molecule inhibitor. Mechanistically, it binds the active,
GTP-bound (ON) state of multiple RAS variants in a tri-complex with the
intracellular chaperone cyclophilin A and thereby suppresses downstream effector
signaling, providing broad pan-KRAS coverage that includes the G12 alleles
(G12D, G12V, and G12R) that predominate in pancreatic ductal adenocarcinoma
(PDAC). This mechanism and pharmacological class are developed in full in the
Introduction (\S 3.1) of this IND and are summarized here only insofar as they
frame the human experience: the drug acts on a validated oncogenic driver in a
disease whose five-year overall survival remains below 13 percent, the third
leading cause of cancer death in the United States and the fourth in the European
Union \cite{pdac060s2030}.

The most significant regulatory milestone in the drug's investigational history to
date is the grant of Breakthrough Therapy designation by the FDA in June 2025 for
previously treated metastatic PDAC harboring KRAS G12 mutations. Breakthrough
Therapy designation is reserved for an investigational drug intended to treat a
serious or life-threatening condition for which preliminary clinical evidence
indicates that the drug may demonstrate substantial improvement over available
therapy on one or more clinically significant endpoints; the designation does not
constitute marketing approval and does not alter the investigational status of the
drug, but it does signal that the FDA has reviewed preliminary clinical data and
found a credible signal of substantial improvement \cite{daraxwhipple,
phase1guidance}. For the purposes of this IND, the designation is relevant as a
qualitative indicator of the maturity and promise of the human experience
accumulated in the metastatic setting, not as a basis for any efficacy claim in the
perioperative, resectable population studied here. Table 9.1 records the marketed and
designation status concisely.

\vspace{0.2em}
\noindent\begin{tabularx}{\textwidth}{L{4.4cm} Y}
\toprule
\textbf{Attribute} & \textbf{Status for daraxonrasib (RMC-6236)} \\
\midrule
Marketing authorization & None in any country; investigational new drug only. \\
Drug class & Oral RAS(ON) multi-selective (pan-RAS) small-molecule inhibitor; binds GTP-bound active mutant RAS in a tri-complex with cyclophilin A. \\
Route and schedule & Oral, once daily. \\
Molecular target population & PDAC with documented KRAS G12 mutations (G12D, G12V, G12R); broad pan-KRAS coverage. \\
FDA Breakthrough Therapy & Granted June 2025 for previously treated metastatic PDAC with KRAS G12 mutations. \\
Marketed safety experience & None; all human safety data derive from the investigational program (\S 9.2). \\
Postmarketing surveillance & Not applicable; no marketed use exists. \\
\bottomrule
\end{tabularx}
\vspace{-0.2cm}
\figcaption{Table 9.1. Marketed and regulatory-designation status of daraxonrasib.
The drug is investigational only; the June 2025 Breakthrough Therapy designation
is a development milestone, not a marketing authorization.}

\subsection{Prior Clinical Research Experience}

The prior clinical research experience with daraxonrasib is concentrated in the
late-phase metastatic-disease program, which comprises two pivotal studies and
supplies the molecular eligibility, dosing, and safety reference framework adapted
for this Phase 1 investigation. The two studies are RASolute 302, a Phase 3 study
in previously treated metastatic PDAC, and RASolve 301, a first-line program; each
is summarized below, followed by the way their experience is translated into the
proposed protocol \cite{daraxwhipple, onpremwhippl}.

\subsubsection*{RASolute 302 (Phase 3, previously treated metastatic PDAC)}

RASolute 302 is the pivotal Phase 3 study of daraxonrasib in patients with
previously treated metastatic PDAC. It establishes the human experience base for
the drug at the dose level that the present protocol adopts as its highest
escalation step: the RASolute 302 recommended Phase 2 dose corresponds to
daraxonrasib 300 mg administered orally once daily, which this IND designates as
dose level 3 (DL3) of its 3+3 escalation, with DL1 at 160 mg and DL2 at 220 mg.
The molecular eligibility of RASolute 302 is broad pan-KRAS, encompassing the G12
alleles, and that breadth is carried forward directly into the present protocol:
eligible participants must have histologically or cytologically confirmed PDAC
with a documented KRAS G12 mutation (G12D, G12V, or G12R), anchoring the present
study's molecular criterion to the validated RASolute 302 eligibility rather than
introducing a narrower or novel mutational requirement \cite{daraxwhipple}. The
safety experience accumulated in RASolute 302, including the characterization of
gastrointestinal effects (nausea, diarrhea, stomatitis), dermatologic and
mucocutaneous events, and hepatic enzyme elevations, provides the reference safety
information against which expectedness is judged for drug-attributed adverse events
in this Phase 1 study, as detailed in the safety framework below.

\subsubsection*{RASolve 301 (first-line program)}

RASolve 301 is the first-line program for daraxonrasib in PDAC and extends the
human experience to the previously untreated metastatic population, including
combination contexts. Its relevance to the present IND is twofold. First, it
broadens the dose and exposure experience that informs the conservative,
sub-RASolute-302 starting dose chosen for this first-in-human-in-context study
(DL1 at 160 mg, well below the 300 mg recommended Phase 2 dose). Second, the
perioperative pause-and-restart advisory specified in this protocol is designed to
respect any first-line combination context that RASolve 301 may establish: the
restart logic keys the resumption day to the postoperative pancreatic fistula
grade and the serum daraxonrasib trough relative to the 0.5 ng/mL threshold, and
the advisory framework is constructed so that it remains valid whether the drug is
given as monotherapy or within a combination regimen \cite{daraxwhipple,
onpremwhippl}.

\subsubsection*{Translation of prior experience into the proposed protocol}

The prior clinical research experience is translated into the present protocol
through three explicit linkages, summarized in Table 9.2. The molecular
eligibility is anchored to the broad pan-KRAS RASolute 302 criteria; the dosing
ladder is anchored to the RASolute 302 recommended Phase 2 dose as its ceiling,
with two lower steps below it; and the reference safety information for drug
causality and expectedness is taken from the metastatic program. The present study
does not seek to reproduce or confirm the efficacy demonstrated in the metastatic
setting; it is a descriptive, non-confirmatory, first-in-human investigation of the
drug in a resectable or borderline-resectable, curative-intent population
undergoing a device-directed pancreaticoduodenectomy, with co-primary endpoints
directed at device and procedure safety, the maximum tolerated dose (MTD) and
recommended Phase 2 dose (RP2D) via 3+3 escalation over a 28-day dose-limiting
toxicity (DLT) window, and technical feasibility \cite{phase1guidance, daraxwhipple}.

\vspace{0.2em}
\noindent\begin{tabularx}{\textwidth}{L{3.6cm} L{4.6cm} Y}
\toprule
\textbf{Prior study} & \textbf{Population and status} & \textbf{Contribution to the present Phase 1 IND} \\
\midrule
RASolute 302 & Phase 3, previously treated metastatic PDAC; broad pan-KRAS eligibility & Defines the molecular eligibility (KRAS G12) used here; its recommended Phase 2 dose (300 mg once daily) is adopted as DL3, the ceiling of the 3+3 ladder; supplies drug reference safety information. \\
RASolve 301 & First-line metastatic PDAC program, including combination contexts & Broadens the dose and exposure experience supporting the conservative 160 mg starting dose; the perioperative pause-and-restart advisory is built to remain valid in any first-line combination context. \\
Breakthrough Therapy (June 2025) & Regulatory designation for previously treated metastatic PDAC with KRAS G12 & Qualitative indicator of the maturity of the metastatic human experience; not a basis for any efficacy claim in the resectable population studied here. \\
\bottomrule
\end{tabularx}
\vspace{-0.2cm}
\figcaption{Table 9.2. How the prior clinical research experience with daraxonrasib
is translated into the eligibility, dosing, and safety reference framework of the
present Phase 1 IND.}

The dose ladder and its 3+3 escalation rule are stated in full in the General
Investigational Plan (\S 4) and are reproduced here only to fix the relationship
between prior experience and the proposed exposures: three participants are enrolled
at each level; if zero of three experience a DLT the cohort escalates, if one of
three does the cohort expands to six, and if one of six does the cohort escalates,
whereas two or more DLTs in six define an exceeded MTD with the MTD declared at the
prior level. The maximum treated sample is 18 participants (three levels times six),
with approximately 36 screened, at a single academic site enrolling approximately
one to two participants per month under staggered sentinel enrollment. Because the
highest planned exposure (DL3, 300 mg once daily) does not exceed the
already-characterized RASolute 302 recommended Phase 2 dose, the prior human
experience directly bounds the upper end of the exposure range proposed here
\cite{daraxwhipple, phase1guidance}.

\subsection{Clinical Care Experience}

The clinical care experience relevant to this IND is the established human
experience with robotic pancreaticoduodenectomy (the Whipple procedure), which
provides the contemporary benchmark against which the device-directed component of
this study will be evaluated, together with the safety-reporting framework that
governs how new human experience is captured once the IND is in effect. The
Whipple procedure is the most technically demanding curative-intent abdominal
oncology operation: it requires resection and reconstruction of four luminal
structures (the pancreas, the duodenum, the common bile duct, and the stomach or
pylorus) and dissection around named vessels (notably the superior mesenteric vein
and the portal vein), and its leading lethal complication is the
fistula-sepsis-hemorrhage cascade graded under the International Study Group of
Pancreatic Surgery (ISGPS) classification of postoperative pancreatic fistula
\cite{onpremwhippl}.

\subsubsection*{The 2025 Dutch nationwide robotic-Whipple benchmark cohort}

The contemporary human baseline for robotic pancreaticoduodenectomy is the 2025
Dutch nationwide cohort of 1000 robotic procedures, which anchors the secondary
endpoint comparisons in this study. That cohort reported a conversion rate of 10.1
percent, Clavien-Dindo grade III or higher complications of 41.3 percent, ISGPS
grade B or C (clinically relevant) postoperative pancreatic fistula of 24.4
percent, and 30-day mortality of 3.9 percent \cite{onpremwhippl}. These four
figures are used descriptively, not as a comparator arm: the present single-arm
study is non-confirmatory, with no power calculation and no formal hypothesis test
against the Dutch cohort, and all comparisons are framed as exact 95 percent
Clopper-Pearson interval estimates around the small-sample proportions observed in
the trial. The benchmark serves to demonstrate that even expert, high-volume
robotic surgery leaves a substantial residual of conversions, high-grade
complications, and clinically relevant fistulae, which are precisely the failure
modes the device-directed component is designed to mitigate through its layered
hard safety limits (the 10 kHz heartbeat bus, the cross-arm emergency stop of 3 ms
or less, per-arm and cumulative tip-force caps, and the five-vessel vascular no-fly
gate) \cite{daraxwhipple}. Table 9.3 places the benchmark alongside the
corresponding study secondary endpoints and the small-sample precision the study
can achieve at its maximum size.

\vspace{0.2em}
\noindent\begin{tabularx}{\textwidth}{L{4.6cm} C{2.6cm} Y}
\toprule
\textbf{Outcome} & \textbf{Dutch 2025 benchmark (n = 1000)} & \textbf{Role in the present study} \\
\midrule
Conversion to open surgery & 10.1\% & Secondary endpoint; device-directed conversion (from malfunction or unsafe behavior) is also a co-primary feasibility input. \\
Clavien-Dindo grade III+ complications (90-day) & 41.3\% & Secondary endpoint; device- or procedure-related serious adverse events through 30 days, including Clavien-Dindo III+, are a co-primary safety endpoint. \\
ISGPS grade B or C pancreatic fistula & 24.4\% & Secondary endpoint and event of special interest; the pancreaticojejunostomy fistula grade is the input to the perioperative drug-restart advisory. \\
30-day mortality & 3.9\% & Secondary endpoint; 30-day and 90-day mortality are both followed. \\
\bottomrule
\end{tabularx}
\vspace{-0.2cm}
\figcaption{Table 9.3. The 2025 Dutch nationwide robotic-Whipple benchmark cohort
(n = 1000) and the corresponding secondary endpoints in the present study. The
benchmark is used descriptively; all study estimates are reported with exact 95
percent Clopper-Pearson confidence intervals.}

Because the study is small by design, its endpoint estimates carry wide exact
confidence intervals, and the benchmark is therefore a contextual reference rather
than a statistical comparator. At the maximum treated sample of 18, an observed
rate of zero of 18 corresponds to 0 percent with a 95 percent Clopper-Pearson
interval of 0 to 18 percent, and one of 18 corresponds to 5.6 percent with an
interval of 0.1 to 27 percent; within a single dose cohort, zero of three
corresponds to an interval of 0 to 71 percent and zero of six to an interval of 0
to 46 percent. These precision figures make explicit that the present study is
designed to characterize the safety and feasibility of the combined intervention,
not to establish superiority or non-inferiority against the Dutch benchmark
\cite{phase1guidance}. The surgical and oncologic care context, including the
duct-to-mucosa pancreaticojejunostomy under controller ring-tension monitoring
(0.30 to 0.60 N), the end-to-side hepaticojejunostomy (0.20 to 0.50 N), and the
antecolic gastrojejunostomy (0.40 to 0.80 N), is specified in full in the protocol
and is referenced here as the established clinical care experience the device is
designed to reproduce within validated envelopes \cite{onpremwhippl}.

\subsubsection*{The safety-reporting framework for ongoing human experience}

The human experience generated by this study will be captured and reported under a
two-stream safety framework that pairs a conventional pharmacovigilance stream for
clinical events with a Physical AI safety stream for device- and AI-related events,
each equally binding. On the clinical stream, any unexpected fatal or
life-threatening suspected adverse reaction is reported to the FDA as soon as
possible but no later than 7 calendar days after initial receipt of the
information, and any other serious and unexpected suspected adverse reaction is
reported no later than 15 calendar days after the determination that the event is
reportable, consistent with 21 CFR \S 312.32 \cite{cfr312paiuni}. A Physical AI
adverse event, per 21 CFR \S 312.32(f), is any untoward occurrence associated with
the operation of the device during the investigation, whether or not caused by it,
including robotic malfunction, AI prediction error, sensor failure, communication
loss, unauthorized system access, digital-twin divergence, and any event requiring
activation of the emergency stop. Six Physical AI reporting triggers are reported
under \S 312.32(g): a serious Physical AI adverse event on the applicable 7-day or
15-day timeline; any system-drug interaction causing incorrect drug administration
regardless of seriousness; any cybersecurity incident within 15 calendar days, or
7 days if it has the potential to cause patient harm; sustained AI model
performance degradation below the pre-specified acceptance criteria for three or
more consecutive procedures or a cumulative 24 hours; any sim-to-real divergence
beyond twice the validated gap tolerance (trajectory deviation greater than 4 mm or
force discrepancy greater than 1.0 N); and any digital-twin discrepancy beyond the
pre-specified accuracy thresholds. For every Physical AI adverse event reported,
the complete hash-chained audit trail covering the period from 24 hours before the
event to 72 hours after the event is preserved under 21 CFR part 11 and made
available to the FDA on request \cite{cfr312paiuni, cfr050paiuni}.

Once the IND is in effect and the 30-day clock has elapsed, post-submission
authoring is driven by this safety data rather than by the initial submission, and
the same generation method that produced the IND maintains the document suite in
internal consistency. Figure 13 shows the post-submission maintenance flow:
serious adverse events and suspected adverse reactions feed the 7-day and 15-day
IND safety reports; emerging pharmacokinetic and pharmacodynamic data and cohort
completion feed the annual Development Safety Update Report (DSUR) and amendments;
data safety monitoring board (DSMB) cohort reviews feed the dose-escalation minutes
and the RP2D memorandum; the principal investigator assesses causality and
scientific justification across all of these; and the LLM then synchronizes the
protocol, consent, and site documents so the whole suite remains internally
consistent \cite{onpremwhippl, ichpaistduni}.

\begin{mermaidfig}
\node[mmgoal] (eff) at (0,3.0) {IND in effect\\(30-day clock elapsed)};
% inputs row
\node[mmin] (sae)  at (-5.2,0.9) {SAE / SUSAR event};
\node[mmin] (pkpd) at (0,0.9)    {Emerging PK / PD,\\cohort completion};
\node[mmin] (dsmb) at (5.2,0.9)  {DSMB cohort review};
% processing row
\node[mmproc] (n7)   at (-6.4,-1.6) {IND safety report\\7-day (fatal /\\life-threatening)};
\node[mmproc] (n15)  at (-3.0,-1.6) {IND safety report\\15-day (other serious,\\unexpected)};
\node[mmproc] (dsur) at (0.4,-1.6)  {Annual DSUR;\\amendments};
\node[mmproc] (min)  at (4.4,-1.6)  {Dose-escalation minutes;\\RP2D memo};
% PI
\node[mmgoal] (pi)   at (0,-3.9) {PI: causality and\\scientific justification};
\node[mmaccent] (sync) at (0,-5.9) {LLM synchronizes protocol\\+ consent + site documents};
\draw[mmedge] (eff) -- (sae);
\draw[mmedge] (eff) -- (pkpd);
\draw[mmedge] (eff) -- (dsmb);
\draw[mmedge] (sae) to[out=-110,in=90,looseness=0.6] (n7.north);
\draw[mmedge] (sae) to[out=-70,in=90,looseness=0.6] (n15.north);
\draw[mmedge] (pkpd) -- (dsur);
\draw[mmedge] (dsmb) to[out=-90,in=90,looseness=0.6] (min.north);
\draw[mmedge] (n7) to[out=-90,in=160,looseness=0.6] (pi.west);
\draw[mmedge] (n15) to[out=-90,in=150,looseness=0.6] (pi.west);
\draw[mmedge] (dsur) -- (pi);
\draw[mmedge] (min) to[out=-90,in=30,looseness=0.6] (pi.east);
\draw[mmedge] (pi) -- (sync);
\end{mermaidfig}
\vspace{-0.4cm}
\figcaption{Figure 13. Post-submission IND maintenance authoring once the IND is in
effect (the 30-day clock elapsed). Safety data drive the authoring: IND safety
reports on the 7-day (fatal or life-threatening) and 15-day (other serious and
unexpected) clocks, the annual DSUR, amendments, and dose-escalation minutes, with
the principal investigator assessing causality and scientific justification and the
LLM synchronizing the protocol, consent, and site documents so the whole suite
stays internally consistent.}

The three-tier oversight that reviews this accumulating human experience comprises
the DSMB (acting at cohort boundaries and on triggered review, with halt
authority), an Independent Safety Monitor (per-procedure, real-time, with stop
authority), and a Physical AI Safety Review Committee (convening at intervals not
exceeding 90 days during active enrollment). The binding halt rules are that one or
more device- or procedure-related serious adverse events within a cohort, or two
device-related deaths at any time, trigger DSMB review and possible halt
\cite{onpremwhippl, congress9510}. This governance ensures that the previous human
experience characterized in this section is continuously updated and acted upon as
the study generates the first direct human experience with the drug in this
device-directed, curative-intent context, the gap that Figure 22 identifies as
first-in-human \cite{clintrustpai}.

\subsection{References}

The references supporting this section are listed in full in the back matter and
are cited inline above. The clinical and computational evidence base is documented
in the daraxonrasib robotic-Whipple and on-premises Whipple sources
\cite{daraxwhipple, onpremwhippl}; the phased-development and Phase 1 expectations
in the Phase 1 guidance \cite{phase1guidance}; the PDAC epidemiology and
metastatic-program context in the disease and program references
\cite{pdac060s2030}; the safety-reporting and Part 11 framework in the 21 CFR
adaptations \cite{cfr312paiuni, cfr050paiuni, ichpaistduni}; the
verification-before-generation and clinical-trust framing in the legislative and
trust references \cite{congress9510, clintrustpai}; and the author's 2025 to 2026
computational program in the digital-twin, quantitative-systems-pharmacology,
in silico Phase 3, and digital-twin-proposal works \cite{fdadigtwinpc, qspmetpancre,
chatgpt100kp, pdacdigtwinp, llmcomp16fda}.
